
\subsubsection{LINGUAGEM DE PROGRAMAÇÃO: \en{KOTLIN}}

Para o desenvolvimento da aplicação móvel, sua construção foi realizada com o \en{software} \en{Android Studio} usando a linguagem de programação \en{Kotlin}. No desenvolvimento de uma aplicação é importante a escolha de uma plataforma, a plataforma escolhida para a aplicação foi o \en{Android}. Este é um sistema operacional para dispositivos móveis, sendo usado em \en{tablets}, celulares, dentre outros. \en{Android Studio} é o ambiente de desenvolvimento oficial de aplicações para a \en{Android}, pode ser usado com  \en{Java} e \en{Kotlin}. \en{Kotlin} é uma linguagem de programação pragmática que combina os paradigmas de Orientação a Objetos e Programação Funcional com foco em interoperabilidade; a interoperabilidade se dá pela qualidade da linguagem de usar códigos que foram escritos em \en{Java} em códigos escritos em \en{Kotlin} \cite{Resende2020KotlinAndroid}. 

Orientação a Objetos é uma paradigma de programação que determina a organização do código em objetos. Um objeto é gerado de uma classe, essa é a especificação de atributos e métodos que o objeto contém em sua definição; cada classe apresenta uma interface consistente para o restante do código sem entregar detalhes inapropriados \cite{GoodrichTamassiaMount2011DataStructuresCpp}. \en{Kotlin} como uma linguagem orientada a objetos adere a esse paradigma.
